\documentclass[10pt,a4paper]{amsart}
\usepackage[margin=19mm]{geometry}
\usepackage{amsmath,amssymb,mathtools}
\usepackage[scr=rsfs,cal=euler]{mathalfa}
\usepackage{xfrac}
\usepackage[unicode,hidelinks,bookmarks=false]{hyperref}
\newcommand{\ie}{i.e.\ }
\newcommand{\cf}{cf.\ }
\newcommand{\resp}{respectively\ }
\newcommand{\Subsection}{\subsection}
\providecommand{\nobreakdash}{\nobreak\hbox{-}}
\setlength{\emergencystretch}{2em}
\multlinegap0pt
\usepackage{amssymb}
\usepackage{enumerate}
\usepackage{xcolor}
\newtheorem{thm}{Theorem}
\title{Integrally Hilbertian rings and \break the polynomial Schinzel hypothesis}
\author{Angelot Behajaina, Pierre D\`ebes, Joachim K\"onig}
\date{Source: arXiv:2502.09959v2 (CC BY 4.0)}
\begin{document}
\maketitle

\noindent\emph{Source and licence.} Integrally Hilbertian rings and \break the polynomial Schinzel hypothesis. Version arXiv:2502.09959v2 (CC BY 4.0), distributed by arXiv under the Creative Commons Attribution 4.0 International licence, http://creativecommons.org/licenses/by/4.0/, as recorded in the arXiv licence metadata for this exact version and shown on the abstract page https://arxiv.org/abs/2502.09959v2. Full licence notice: \texttt{licenses/arxiv-2502.09959-CC-BY-4.0.txt}.\par\smallskip

\noindent\textit{Editorial scope.} Counted unit: the article's thm thm:main3 (source0.tex lines 2017--2102). The article's complete original proof of it is delivered with it (source0.tex lines 2457--2574, one \texttt{proof} environment of 118 lines, which the article places away from the statement). No mathematics is written, repaired or re-derived: the statement and the proof are the article's own lines, in the article's own notation, and no retained line is substituted or re-flowed.  No further source-local context is needed: every cross-reference of every retained line resolves inside the two delivered spans.  Nothing was omitted from any retained span, so the package carries no omission record.  No retained line contains a citation, so the wrapper carries no bibliography and no bibliography entry.  No retained line contains a figure environment, an image-inclusion command or drawing code, so no image is delivered and the package declares no asset dependency.  Editorial formatting only, in addition: an AMS \texttt{amsart} preamble that restates the article's own declarations and macros used by the retained lines, namely the environments \texttt{thm} on separate counters, together with the standard packages those lines need.  The retained environments print in this document as Theorem 1 in order of appearance, whereas the article prints the same items with the numbering of its own full document.  The complete native source of the article as distributed in the arXiv e-print for this exact version is bundled unchanged as \texttt{sources/173290/source0.tex}.
\begin{thm} \label{thm:main3}
Assume that ${\mathcal Z}$ is an integrally Hilbertian ring.
Let $P_1,\ldots,P_s \in {\mathcal Z}[\underline{T},\underline{Y}]$ be $s\geq 1$ polynomials, irreducible in $\mathcal{Q}[\underline{T},\underline{Y}]$, and such that the product $P_1\cdots P_s$ has no nonunit divisor in ${\mathcal Z}$.
For each $i=1,\ldots,k$, let ${\underline d}_i=(d_{i1},\ldots,d_{in})$ be an $n$-tuple of nonnegative integers. Assume that \vskip 1mm

\noindent
{\rm (1)} $\deg_{\underline Y}(P_l) \geq 1$ for $l=1,\ldots,s$ or ${\underline d}_i \not= (0,\ldots,0)$ for $i=1,\ldots,k$. 
\vskip 2mm 

\noindent
Consider the subset
\vskip 1mm

\centerline{${\mathcal S}_{{\mathcal Z}[\underline Y]}(\underline P,{\underline d}_1,\ldots,{\underline d}_n) \subset {\mathcal P}ol_{{\mathcal Z},n,\underline{d}_1} \times \cdots \times {\mathcal P}ol_{{\mathcal Z},n,\underline{d}_k}$} 
\vskip 2mm

\noindent
consisting of all $k$-tuples $(M_1(\underline Y), \ldots, M_k(\underline Y))$ such that   \vskip 1,5mm

\noindent
{\rm (2)} \hskip 10mm $P_i\hskip 2pt (M_1(\underline Y), \ldots, M_k(\underline Y),Y_1, \ldots, Y_n)$
is irreducible in ${\mathcal Z}[\underline Y]$, $i=1,\ldots, s$.
%\vskip 0,5mm

%\hskip 12mm and $\deg_{Y_j}(M_i) = d_{ij}$, $j=1,\ldots,n$, $i=1,\ldots, k$.
\vskip 2mm

\noindent
Then we have the following three conclusions:
\vskip 1mm


\noindent
{\rm (a)} Assume further that
\vskip 0,5mm

\noindent
{\rm (3.a)} \hskip 25mm $\ell({\underline d}_i)> \sum_{j=1}^s \deg_{T_i}(P_j)$, \hskip 3mm $i=1,\ldots,k$.
\vskip 0,5mm

\noindent
Then the subset ${\mathcal S}_{{\mathcal Z}[\underline Y]}(\underline P,{\underline d}_1,\ldots,{\underline d}_n)$ is Zariski-dense in ${\mathcal P}ol_{{\mathcal Z},n,\underline{d}_1} \times \cdots \times {\mathcal P}ol_{{\mathcal Z},n,\underline{d}_k}$.


%it is possible to substitute some polynomials $M_i(\underline Y)\in {\mathcal Z}[\underline Y]$ of degree $d_{ij}$ in $Y_j$ ($j=1,\ldots,n$) for each parameter $T_i$ ($i=1,\ldots,k$) in the original polynomials $P_1,\ldots,P_s$ in such a way that 


%\noindent
%Furthermore, we have these adjusted conclusions in the following special cases:
\vskip 2mm

\noindent
{\rm (b)} If ${\mathcal Z}$ is a {\it near UFD}, then the same holds with {\rm (3.a)}
%condition on $\ell(\underline d_1), \ldots, \ell(\underline d_k)$ 
replaced by 
\vskip 0,5mm

\noindent
{\rm (3.b)} \hskip 20mm $2^{\ell({\underline d}_i)} > \sum_{j=1}^s \deg_{T_i}(P_j)$, \hskip 3mm $i=1,\ldots,k$.
\vskip 2mm

\noindent
{\rm (c)} If $k=1$ and $\underline {d}_1 = \underline d = (d_1,\ldots,d_n)$, 
the same holds with {\rm (3.a)}
replaced by 
%the conclusion holds under this condition on $\underline {d}_1 = (d_1,\ldots,d_n)$:
\vskip 1mm

\noindent
{\rm (3.c)} \hskip 20mm $\sum_{j=1}^n d_j > \sum_{i=1}^s\deg_{\underline{Y}}(P_i)$.
\vskip 0,7mm
\noindent
and, if additionally ${\mathcal Z}$ is a {near UFD}, by 
\vskip 1mm


\noindent 
{\rm (3.c')}$ \hskip 20mm\sum_{j=1}^n d_j  > \max_{1 \leq i \leq s}\deg_{\underline{Y}}(P_i)$.

%\centerline{$\prod_{j=1}^n (d_j+1)> \sum_{i=1}^s \deg_{T}(P_i)$\hskip 5mm  or \hskip 5mm $\sum_{j=1}^n d_j \geq \sum_{i=1}^s\deg_{\underline{Y}}(P_i)+1$.}
%\vskip 2mm
%\noindent
%and, if ${\mathcal Z}$ is a {near UFD}, by
%\vskip 2mm
%\centerline{$2^{\prod_{j=1}^n (d_j+1)} > \sum_{i=1}^s \deg_{T}(P_i)$ \hskip 5mm or \hskip 5mm $\sum_{j=1}^n d_j  \geq \max_{1 \leq i \leq s}\deg_{\underline{Y}}(P_i)+1$.}
\end{thm}

\begin{proof} Let $p$ be a nonunit of ${\mathcal Z}$, $p\not=0$. By assumption, the polynomial $\Pi(\underline T,\underline Y)$ is nonzero in $({\mathcal Z}/p{\mathcal Z})[\underline T,\underline Y]$. We start with case (b) of Theorem \ref{thm:main3}.
\vskip 1mm

\noindent
\emph{Case {\rm (b)}: ${\mathcal Z}$ is a near UFD and {\rm (3.b)} holds: $2^{\ell({\underline d}_i)} > \sum_{j=1}^s \deg_{T_i}(P_j)$} ($i=1,\ldots,k$). As ${\mathcal Z}$ is a near UFD, one may assume that $p$ is a prime. For each $i=1,\ldots,k$, the subset
\vskip 2mm

\centerline{$\{M_{\underline {Q_i}}({\overline{\underline{\theta_i}}}, \underline Y) \hskip 2pt | \hskip 2pt \overline{\underline{\theta_i}} \in ({\mathcal Z}/p{\mathcal Z})^{\ell_i+1}\} \subset ({\mathcal Z}/p{\mathcal Z})[\underline Y]$,}
\vskip 2mm

\noindent
is of cardinality $|{\mathcal Z}/p{\mathcal Z}|^{\ell_i+1} \geq 2^{\ell_i+1}$. 
Assumpiion (3.b) rewrites as $2^{\ell_i+1}> \deg_{T_i}(\Pi)$, $i=1,\ldots,k$. Since ${\mathcal Z}/p{\mathcal Z}$ is a domain, it follows that there exists a $k$-tuple 
\vskip 1mm

\centerline{$\underline M_{\underline Q}( \overline{\underline{\Theta}},\underline Y)=(M_{\underline {Q_1}}({\overline{\underline{\theta_1}}}, \underline Y), \ldots, M_{\underline {Q_k}}({\overline{\underline{\theta_k}}}, \underline Y)) \in (({\mathcal Z}/p{\mathcal Z})[\underline Y])^k$}
\vskip 1mm

\noindent
such that
$\Pi\hskip 1pt (\underline M_{\underline Q}(\overline{\underline{\Theta}},\underline Y),\underline Y)$ is nonzero in $({\mathcal Z}/p{\mathcal Z})[\underline Y]$.
 Lifting the $\overline{\underline{\theta_i}}$ to some elements of ${\mathcal Z}$ provides a point ${\underline{\Theta}} = (\underline {\theta_1},\ldots,\underline{\theta_k})\in \prod_{i=1}^k {\mathcal Z}^{\ell_i+1}$ such that 
 \vskip 1mm
 
 \centerline{$\underline M_{\underline Q}( \underline{\Theta},\underline Y)=(M_{\underline {Q_1}}({\underline{\theta_1}}, \underline Y), \ldots, M_{\underline {Q_k}}(\underline{\theta_k}, \underline Y)) \in ({\mathcal Z}[\underline Y])^k$}
 
 \vskip 1,5mm

 \noindent
  has the property that $\Pi\hskip 1pt (\underline M_{\underline Q}(\underline{\Theta},\underline Y),\underline Y)$ is not divisible by $p$.  
Conclude that $\underline{\Theta}$ is a special value of $\underline{\Lambda}$ that makes the polynomial $\Pi\hskip 1pt (\underline M_{\underline Q}(\underline{\Lambda},\underline Y),\underline Y)$ nonzero modulo $p$, \hbox{i.e.} $p$ is not a fixed divisor of $F_1\cdots F_s$ \hbox{w.r.t.} $\underline{\Lambda}$.
\vskip 2mm

\noindent
\emph{Case {\rm (a)}: ${\mathcal Z}$ is  integrally Hilbertian and \hbox{\rm (3.a)} holds: $\ell({\underline d}_i) > \sum_{j=1}^s \deg_{T_i}(P_j)$ ($i=1,\ldots,k$)}. The ring ${\mathcal Z}/p{\mathcal Z}$ is no longer integral in general. We will however be able to adjust the argument by working with a smaller subset of tuples of polynomials $\underline M_{\underline Q}(\underline{\Theta},\underline Y)$
than the one used in Case (b).

Specifically, for each $i=1,\ldots,k$, consider the following subset:
\vskip 1mm

\centerline{${\mathcal S}_i=\{ Q_{il} \hskip 2pt | \hskip 2pt  l=0,\ldots,\ell_i\} \subset {\mathcal Z}[\underline Y]$.}

\vskip 1mm

\noindent
We claim that the corresponding elements, regarded in $({\mathcal Z}/p{\mathcal Z})[\underline Y]$, have this property: any difference between two distinct such elements is not a zero divisor in $({\mathcal Z}/p{\mathcal Z})[\underline Y]$. Namely, if $P\in ({\mathcal Z}/p{\mathcal Z})[\underline Y]$ is nonzero and $Q$ is the largest monomial in $P$ for the lexicographical order, then the largest monomial in  the polynomial $Q_{il} P$ is $Q_{il} Q$. For two different $Q_{il}, Q_{i l^\prime} \in {\mathcal S}_i$, we have $Q_{il} Q\not= Q_{i l^\prime} Q$ and so $Q_{il} P\not= Q_{i l^\prime} P$, which proves the claim. Deduce next that the product of all nonzero differences $(Q_{il}- Q_{i l^\prime})$ with $Q_{il}, Q_{i l^\prime} \in {\mathcal S}_i$ is not a zero divisor in $({\mathcal Z}/p{\mathcal Z})[\underline Y]$.


 It follows from this conclusion that a nonzero polynomial in one variable of degree $\leq \ell_i$ and with coefficients in $({\mathcal Z}/p{\mathcal Z})[\underline Y]$ cannot vanish at every one of the $\ell_i+1$ elements $Q_{il}$ with $ l=0,\ldots,\ell_i$. Indeed the determinant of the linear system corresponding to the $\ell_i+1$ vanishing conditions -- a Van Der Monde determinant -- is precisely the product of differences in $({\mathcal Z}/p{\mathcal Z})[\underline Y]$ considered above. By construction it is not a zero divisor in $({\mathcal Z}/p{\mathcal Z})[\underline Y]$, and so only the zero polynomial could be a  solution to the system, which is excluded.

 
Conclude: as $\ell_i=\ell(\underline d_i) - 1 \geq \deg_{T_i}(\Pi)$ for each $i=1,\ldots,k$, there exist $l_i\in \{0,\ldots,\ell_i\}$ ($i=1,\ldots,k$) such that the $k$-tuple $\underline Q_{\underline l}=(Q_{1l_1},\ldots,Q_{kl_k}) \in ({\mathcal Z}[\underline Y])^k$ satisfies the following.
\vskip 1mm

\centerline{$\Pi\hskip 1pt (\underline Q_{\underline l},\underline Y)$ is not divisible by $p$.}
\vskip 1mm

\noindent
Conclude as in case (b) that $p$ is not a fixed divisor of $F_1\cdots F_s$ \hbox{w.r.t.} $\underline{\Lambda}$.
\vskip 2mm

\noindent
\emph{Case {\rm (c)}: ${\mathcal Z}$ integrally Hilbertian, $k=1$ and \hbox{\rm (3.c)} holds
$\sum_{j=1}^n d_j > \sum_{i=1}^s \deg_{\underline Y}(P_i)$.}
Write $T$ for $T_1$, $\underline \lambda$ for $\underline{\lambda_1}$ and $\ell$ for $\ell(\underline d)-1$. We may assume that $r=\deg_T(\Pi) \geq 1$. Set
\vskip 1mm

\centerline{$\Pi= P_1 \dots P_s=Z_0(\underline{Y})+Z_1(\underline{Y})T+\dots+Z_r(\underline{Y})T^r$.}
\vskip 1mm

\noindent
Assume on the contrary that $F_1\cdots F_s$ has a fixed divisor \hbox{w.r.t.} $\underline{\lambda}$, \hbox{i.e.}, there is a nonunit $p\in \mathcal{Z}$, $p\not=0$, such that 
\vskip 1mm

\centerline{$(F_1\cdots F_s)(\underline{m},\underline{Y})\equiv 0 \pmod{p}$\hskip 2mm for all $\underline{m} \in \mathcal{Z}^{\ell+1}$.}
\vskip 1mm


\noindent
In particular, for $\underline m$ corresponding to the monomial $Y_1^{d_1}\dots Y_n^{d_n}$, we obtain:
\vskip 1,5mm

\centerline{$Z_0(\underline{Y})+Z_1(\underline{Y})Y_1^{d_1}\dots Y_n^{d_n}+Z_2(\underline{Y})Y_1^{2d_1}\dots Y_n^{2d_n}+\dots+Z_r(\underline{Y})Y_1^{r d_1}\dots Y_n^{r d_n}\equiv 0 \pmod{p}$.}
\vskip 1,5mm


%Consider the unique $\underline{m}_0 \in \mathcal{Z}^{N_{\underline{d}}}$ such that $M_{\underline{d}}(\underline{m}_0,\underline{Y})=Y_1^{d_1}\dots Y_n^{d_n}$. Then

%\begin{equation}\label{eq:condgmo1}
%Z_0(\underline{Y})+Z_1(\underline{Y})Y_1^{d_1}\dots Y_n^{d_n}+Z_2(\underline{Y})Y_1^{2d_1}\dots Y_n^{2d_n}+\dots+Z_r(\underline{Y})Y_1^{r d_1}\dots Y_n^{r d_n}\equiv 0 \pmod{p}.
%\end{equation}


\noindent
We claim that for every $0 \leq i<j \leq r$ such that $Z_i(\underline{Y})$ and $Z_j(\underline{Y})$ are nonzero, the nonzero monomials appearing in $Z_i(\underline{Y})Y_1^{id_1}\dots Y_n^{id_n}$ are different from the ones in $Z_j(\underline{Y})Y_1^{jd_1}\dots Y_n^{j d_n}$. Indeed for nonzero monomials $A$ and $B$ in $Z_i(\underline{Y})Y_1^{id_1}\dots Y_n^{id_n}$ and $Z_j(\underline{Y})Y_1^{jd_1}\dots Y_n^{j d_n}$ respectively, we have:
\begin{align*}
\deg_{\underline{Y}}(A) &\leq \deg_{\underline{Y}}(Z_i)+i\left(\sum_{i=1}^n d_i\right)\\
&\leq \deg_{\underline{Y}}(\Pi)+i\left(\sum_{i=1}^n d_i\right)\\
&= \sum_{i=1}^s\deg_{\underline{Y}}(P_i)+i\left(\sum_{i=1}^n d_i\right)\\
&<\left(\sum_{i=1}^n d_i\right)+i\left(\sum_{i=1}^n d_i\right) \quad \textrm{(by assumption)}\\
&\leq j\left(\sum_{i=1}^n d_i\right)\leq \deg_{\underline{Y}}(B).
\end{align*}
Thus $A \neq B$.
Hence, for each $i=1,\ldots, r$, we have $Z_i(\underline{Y})Y_1^{i\cdot d_1}\dots Y_n^{i \cdot d_n}\equiv 0 \pmod{p}$, and so $Z_i(\underline{Y}) \equiv 0 \pmod{p}$. Therefore $p$ divides $\Pi$, a contradiction.

%Therefore, $P_1\dots P_s$ is not primitive \hbox{w.r.t.} $\mathcal{Z}$, a contradiction. Consequently, $G(\underline{\lambda},\underline{Y})$ has no fixed divisor w.r.t. $\underline{\lambda}$ in ${\rm Nonunit}\hskip 1pt\mathcal{Z}$.
\vskip 2mm


Finally, assume as in the special case of Theorem \ref{thm:main3}(c) that \emph{$\mathcal{Z}$ is a near UFD and that \hbox{\rm (3.c')} holds: $\sum_{j=1}^nd_j> \max_{1 \leq i \leq s}\deg_{\underline{Y}}(P_i)$.} If $p$ is a fixed divisor of $F_1\cdots F_s$ \hbox{w.r.t.} $\underline \lambda$, then $\Pi(Y_1^{d_1}\dots Y_n^{d_n},\underline{Y})\equiv 0 \pmod{p}$. Without loss of generality, one may assume that $p$ is prime. Since $(\mathcal{Z}/p\mathcal{Z})[\underline{Y}]$ is an integral domain, there exists $1 \leq j \leq s$ such that 
\vskip 1mm

\centerline{$P_{j}(Y_1^{d_1}\dots Y_n^{d_n},\underline{Y})\equiv 0 \pmod{p}$.}
\vskip 1mm

\noindent
In the second part of the argument above, with $P_j$ replacing $P_1\cdots P_s$, we conclude that $P_j \equiv 0 \pmod{p}$, a contradiction.
\end{proof}

\end{document}
